\begin{figure*}[!htb]
\small
\begin{tcolorbox}[
    colback=white, 
    colframe=blue3,
    boxrule=0.5mm,
    arc=0mm,
    outer arc=0mm,
    title=Limited Functionality Tools: Optimal Tool Selection - Generation Prompt
]
Task Description:\\
You need to go ahead and generate samples of tool usage with as much diversity as possible, based on the current dimension, the requirements requirements of the task, and the sample references provided.\\

Current dimension:\\
Multiple Appropriate Tools Dimension - Optimal Tool Selection\\

Task Requirement:\\
- Generate new samples taking into account as much as possible the diversity of user query and available tools.
- First, make sure that the task description is a goal that requires the use of multiple tools to accomplish, and provide the corresponding tools in the list of available tools using numerical numbers.\\
- Use "// Raw Input //" and "// Raw output //" to label the inputs and outputs of the samples.\\
- The samples need to consider a realistic scenario as much as possible. Multiple tools need to be available in the list of available tools in the raw input. Note that there will be multiple suitable tools that all take on the same or very similar functions in the current dimension, e.g., for a language detection task, tool A may have a higher recognition rate for some minor languages, while tool B is more accurate on common languages. One needs to be weighed against the specific needs of the task (whether comprehensiveness is more important or accuracy for a particular language is more important); for a text sentiment analysis task, tool C can provide simple keyword-based sentiment analysis, while tool D is capable of deeper sentiment understanding and contextual analysis. In this case, if the user needs a quick and shallow analysis, tool C may be preferred; if deeper sentiment understanding is needed, tool D may be preferred. One needs to choose the right tool based on the specific situation, e.g., the complexity of the text and the detailed needs of the user; and so on.\\

At build time, you need to:\\
First give a list of overall tasks and then express specific task requirements separately. In the list of available tools, make sure to provide tools that can fulfill these requirements, while adding a small number of tools that do not directly help the task to ensure the complexity of the task and the diversity of tools. At the same time, in the current dimension, you need to build multiple tools with the same functionality but different focuses, and emphasize a particular focus in the task description to achieve the effect of selecting the most effective tool. Finally, according to these requirements, complete the corresponding task planning process in the output, keeping in mind that each step corresponds to a different tool.\\

Sample Reference:\\
==== Reference Sample Start ====\\
// Raw Input //\\
\colorbox{blue3!50}{\texttt{[sample\_input]}}\\

// Raw output //\\
\colorbox{blue3!50}{\texttt{[sample\_outputl]}}\\
==== Reference Sample End ====\\
Now, please combine the reference sample and the corresponding build steps, try to select tasks from different domains and scenarios, and use different types of tools to make the new sample different from the reference sample in terms of task content and context. We start generating new single samples corresponding to the dimensions.\\
Multiple Appropriate Tools Dimension - Optimal Tool Selection generation sample (one):\\
==== Sample Start ====
\end{tcolorbox}
\caption{Generation prompt for limited functionality tools scenario, optimal tool selection subtask.}
\label{fig:best-prompt}
\end{figure*}